\documentclass[11pt]{article}
\usepackage[T1]{fontenc}
\usepackage{lmodern}
\usepackage{amsmath,amssymb,amsthm}
\usepackage{graphicx}
\usepackage[margin=1in]{geometry}
\usepackage[colorlinks=true,linkcolor=blue,citecolor=blue,urlcolor=blue]{hyperref}

\newtheorem{theorem}{Theorem}
\newtheorem{lemma}{Lemma}
\newcommand{\Dhet}{\mathcal D_{\rho(\cdot)}}
\newcommand{\Qhet}{\mathcal Q_{\rho(\cdot)}}
\newcommand{\sphere}{\mathbb S^{n-1}}

\title{Kernel Triviality for Every Admissible\\
Heterogeneous Horizon}
\author{Run \texttt{fa\_banach\_001}, agent \texttt{agent\_lane\_01}}
\date{July 23, 2026}

\begin{document}
\maketitle

\section*{Source question}

Kreisbeck and Schönberger introduce heterogeneous nonlocal gradients whose
interaction horizon \(\delta(x)\) stays inside the domain and vanishes
smoothly at its boundary \cite{KS}. Their Theorem~4.17 proves, with no
amplitude restriction, that the nullspace is finite-dimensional and consists
of smooth functions. When the maximal horizon is sufficiently small
(``mild variation''), they prove that the nullspace consists exactly of
constants. Remark~4.18 asks whether the same conclusion holds without that
smallness assumption.

\begin{center}
\includegraphics[width=\linewidth]{figures/open_problem_crop.png}
\end{center}

The theorem below gives an affirmative answer in every space dimension.

\section*{Setting}

Let \(\Omega\subset\mathbb R^n\) be a bounded connected Lipschitz domain and
let \(p\in(1,\infty)\). The horizon \(\delta\in C^\infty(\mathbb R^n)\)
satisfies the source paper's assumptions; in particular,
\[
 \delta>0\ \hbox{in }\Omega,\qquad
 \delta=0\ \hbox{on }\Omega^c,\qquad
 \delta(x)\leq\operatorname{dist}(x,\Omega^c).
\]
Let \(\bar\rho_1\) be the radial representative of the source kernel and put
\[
 Q(r)=Q_{\rho_1}(r)=\int_r^1\frac{\bar\rho_1(s)}s\,ds,
 \qquad
 g(r)=r^{n-2}\bar\rho_1(r).
\]
Hypothesis H1 says that \(g\) is decreasing. The averaging operator on
\(\mathbb R^n\) is
\[
 (\Qhet f)(x)=
 \begin{cases}
 \displaystyle
 \int_{\mathbb R^n}\delta(x)^{-n}
 Q\!\left(\frac{|x-y|}{\delta(x)}\right)f(y)\,dy,
 &x\in\Omega,\\[1.2ex]
 f(x),&x\in\Omega^c.
 \end{cases}
\]
The source normalization gives \(\int_{\mathbb R^n}Q(|z|)\,dz=1\), and
\[
 \Dhet u=\Qhet^\Omega\nabla u.
\]

\section*{Idea of the proof}

The source proof already turns a kernel element \(u\) into a smooth,
compactly supported vector field \(\widetilde w\) with
\(\Qhet\widetilde w=0\) and
\(\widetilde w|_\Omega=\nabla u\). The remaining task is therefore a scalar
injectivity statement for the components of \(\widetilde w\).

For \(n\geq2\), solve \(\Delta F=\widetilde w_j\) by a Newtonian potential.
Two radial integrations by parts show that \(\Qhet\Delta F\) is a
nonnegative mixture of
\[
 \text{spherical mean of \(F\) about \(x\)}-F(x).
\]
This gives a maximum principle. The Newtonian potential decays at infinity,
so it must vanish. In dimension two, decay requires zero total mass; that
mass vanishes because \(\widetilde w=\nabla u\) in \(\Omega\), while \(u\)
is constant on every boundary curve. Dimension one is handled by the
corresponding centered-second-difference maximum principle.

\section*{The spherical-mean identity}

\begin{lemma}\label{lem:spherical}
Let \(F\in C^2(\mathbb R^n)\), \(x\in\Omega\), and \(d=\delta(x)>0\). Define
\[
 M_x(r)=\frac1{|\sphere|}
 \int_{\sphere}F(x+dr\theta)\,d\theta,\qquad 0\leq r\leq1.
\]
Then
\begin{equation}\label{eq:spherical}
 (\Qhet\Delta F)(x)
 =\frac{|\sphere|}{d^2}
 \int_0^1(-g'(r))\bigl(M_x(r)-M_x(0)\bigr)\,dr.
\end{equation}
The integral is understood improperly at zero.
\end{lemma}

\begin{proof}
The standard spherical-mean identity gives
\[
 M_x''(r)+\frac{n-1}{r}M_x'(r)
 =\frac{d^2}{|\sphere|}
 \int_{\sphere}\Delta F(x+dr\theta)\,d\theta.
\]
Using polar coordinates and then integrating by parts,
\begin{align*}
 (\Qhet\Delta F)(x)
 &=\frac{|\sphere|}{d^2}
   \int_0^1Q(r)\bigl(r^{n-1}M_x'(r)\bigr)'\,dr\\
 &=\frac{|\sphere|}{d^2}
   \int_0^1r^{n-2}\bar\rho_1(r)M_x'(r)\,dr\\
 &=\frac{|\sphere|}{d^2}
   \int_0^1(-g'(r))\bigl(M_x(r)-M_x(0)\bigr)\,dr.
\end{align*}
Here \(Q'=-\bar\rho_1/r\).

We verify the singular boundary terms. At \(r=1\), both \(Q\) and \(g\)
vanish because \(\rho_1\) has compact support in the open unit ball. At zero,
\(M_x'(r)=O(r)\) and \(M_x(r)-M_x(0)=O(r^2)\). Moreover,
\[
 r^nQ(r)
 =\int_r^1\left(\frac rs\right)^n
     s^{n-1}\bar\rho_1(s)\,ds\longrightarrow0
\]
by dominated convergence and radial integrability. Finally, monotonicity of
\(g\) gives
\[
 \int_{r/2}^{r}s^{n-1}\bar\rho_1(s)\,ds
 =\int_{r/2}^{r}s\,g(s)\,ds
 \geq\frac38r^2g(r)
 =\frac38r^n\bar\rho_1(r).
\]
The left side tends to zero. Thus
\(Q(r)r^{n-1}M_x'(r)\to0\) and
\(g(r)(M_x(r)-M_x(0))\to0\), completing both integrations by parts.
\end{proof}

\section*{A maximum principle}

Define
\[
 \mathcal L F:=\Qhet\Delta F.
\]

\begin{lemma}\label{lem:max}
Let \(F\in C^2(\mathbb R^n)\) satisfy \(\mathcal LF=0\). If
\(\overline\Omega\subset B_R\), then
\[
 \inf_{\partial B_R}F\leq F(x)\leq\sup_{\partial B_R}F
 \qquad(x\in B_R).
\]
\end{lemma}

\begin{proof}
Since \(g\) is decreasing, \(-g'\geq0\). Lemma~\ref{lem:spherical} therefore
implies that \(\mathcal LG(x_0)\leq0\) whenever \(x_0\in\Omega\) is a maximum
of \(G\) over \(B_R\): the interaction ball is contained in
\(\Omega\subset B_R\), and every spherical mean is at most \(G(x_0)\).
If \(x_0\in\Omega^c\), then \(\mathcal LG(x_0)=\Delta G(x_0)\leq0\) at an
interior maximum.

For \(\varepsilon>0\), put \(G_\varepsilon=F+\varepsilon|x|^2\). The
normalization of \(Q\) gives
\[
 \mathcal L G_\varepsilon=2n\varepsilon>0.
\]
Consequently \(G_\varepsilon\) cannot attain its maximum in the interior of
\(B_R\). Hence, for \(x\in B_R\),
\[
 F(x)+\varepsilon|x|^2
 \leq\sup_{\partial B_R}F+\varepsilon R^2.
\]
Letting \(\varepsilon\downarrow0\) proves the upper bound. Applying the same
argument to \(F-\varepsilon|x|^2\) and interior minima proves the lower bound.
\end{proof}

\begin{lemma}\label{lem:scalar}
Let \(f\in C_c^\infty(\mathbb R^n)\) and suppose \(\Qhet f=0\).
\begin{enumerate}
\item If \(n\geq3\), then \(f=0\).
\item If \(n=2\) and \(\int_{\mathbb R^2}f=0\), then \(f=0\).
\end{enumerate}
\end{lemma}

\begin{proof}
Let \(F=\Gamma_n*f\), where \(\Gamma_n\) is the Newtonian kernel normalized
so that \(\Delta\Gamma_n=\delta_0\); see, for example, \cite[Section~2.2]{Evans}.
Then \(F\in C^\infty(\mathbb R^n)\), \(\Delta F=f\), and
\(\mathcal LF=\Qhet f=0\).

For \(n\geq3\), compact support of \(f\) gives
\(F(x)=O(|x|^{2-n})\), so \(F(x)\to0\) at infinity. For \(n=2\),
\[
 F(x)=\frac1{2\pi}\int_{\mathbb R^2}\log|x-y|\,f(y)\,dy.
\]
The coefficient of \(\log|x|\) is \((2\pi)^{-1}\int f\), hence the stated
zero-mass assumption gives \(F(x)=O(|x|^{-1})\). In either case,
\(F\to0\) uniformly on large spheres.

Apply Lemma~\ref{lem:max} on \(B_R\) and let \(R\to\infty\). It follows that
\(F=0\), and therefore \(f=\Delta F=0\).
\end{proof}

\section*{The two-dimensional moment}

\begin{lemma}\label{lem:moment}
Suppose \(n=2\), \(u\in C^\infty(\overline\Omega)\), and
\(w\in C_0^\infty(\overline\Omega;\mathbb R^2)\) satisfies
\(w=\nabla u\) in \(\Omega\). Then
\[
 \int_{\mathbb R^2}w_j\,dx=0,\qquad j=1,2.
\]
\end{lemma}

\begin{proof}
Because the zero extension of \(w\) is smooth, \(w=0\) on
\(\partial\Omega\). Hence \(\nabla u=0\) there. The boundary of a bounded
planar Lipschitz domain is a finite union of closed Lipschitz curves
\(\Gamma_\ell\), and \(u\) is constant, say \(c_\ell\), on each curve.
The divergence theorem gives
\[
 \int_{\mathbb R^2}w_j\,dx
 =\int_\Omega\partial_j u\,dx
 =\sum_\ell c_\ell
   \int_{\Gamma_\ell}\nu_j\,d\mathcal H^1.
\]
The integral of the normal vector around every closed Lipschitz curve is
zero (equivalently, apply the divergence theorem to a constant vector field
on the region bounded by that curve). Every summand vanishes.
\end{proof}

\section*{Dimension one}

\begin{lemma}\label{lem:one}
In the source setting with \(n=1\), \(\Dhet u=0\) implies that \(u\) is
constant, for every admissible horizon.
\end{lemma}

\begin{proof}
The no-smallness part of the source theorem gives
\(u\in C^\infty([a,b])\), where \(\Omega=(a,b)\). Choose \(F\) with
\(F''=u'\), put
\[
 h(r)=\frac{\bar\rho_1(r)}r,
\]
and note that H1 makes \(h\) decreasing and \(-h'>0\) near zero. The identity
\(\Dhet u=\Qhet^\Omega u'\), followed by two one-dimensional integrations by
parts, gives
\[
0=\frac1{\delta(x)^2}\int_0^1(-h'(t))
\bigl(F(x+\delta(x)t)+F(x-\delta(x)t)-2F(x)\bigr)\,dt.
\]
Indeed, H1 gives an \(\iota>0\) such that \(t^\iota h(t)\) is decreasing
near zero, while H0 gives \(h(t)>0\) there; hence
\[
 -h'(t)\geq\frac{\iota}{t}h(t)>0.
\]
The endpoint terms vanish by the same estimates used in
Lemma~\ref{lem:spherical}.

Subtract from \(F\) its secant line on \([a,b]\), obtaining \(G\) with
\(G(a)=G(b)=0\). At a positive interior maximum every centered second
difference is nonpositive. Since \(-h'>0\) near zero, equality forces \(G\)
to equal its maximum in a neighborhood. Thus the maximum set is nonempty,
open, and closed in \((a,b)\), contradicting the boundary values. The same
argument excludes a negative minimum. Hence \(G=0\), \(F\) is affine, and
\(u'=F''=0\).
\end{proof}

\section*{Full solution}

\begin{theorem}\label{thm:main}
Under all the hypotheses of Kreisbeck and Schönberger
\cite[Theorem~4.17]{KS}, but with no mild-variation or small-amplitude
assumption on the admissible horizon, one has
\[
 \ker\Dhet
 =\{u:\Omega\to\mathbb R:\ u\text{ is constant}\}.
\]
This holds for every \(n\geq1\) and every \(p\in(1,\infty)\).
\end{theorem}

\begin{proof}
Constants plainly lie in the kernel. Let \(u\in\ker\Dhet\). The
no-smallness part of the source proof constructs
\[
 \widetilde w
 =\nabla\widetilde u-\mathcal D_{\rho(\cdot)}^{\mathbb R^n}\widetilde u
 \in C_0^\infty(\overline\Omega;\mathbb R^n)
\]
such that
\[
 \Qhet\widetilde w=0,\qquad
 \widetilde w|_\Omega=\nabla u.
\]
These are equations (4.15)--(4.16) and the surrounding argument in
\cite[Theorem~4.17]{KS}; neither uses mild variation.

If \(n=1\), Lemma~\ref{lem:one} gives the conclusion directly. Suppose
\(n\geq2\). Apply Lemma~\ref{lem:scalar} to each component
\(\widetilde w_j\). For \(n\geq3\), no additional condition is needed. For
\(n=2\), Lemma~\ref{lem:moment} supplies
\(\int\widetilde w_j=0\). Thus \(\widetilde w=0\) in every case. Consequently
\(\nabla u=\widetilde w|_\Omega=0\), and connectedness of \(\Omega\) makes
\(u\) constant.
\end{proof}

\section*{Verification notes}

The script \texttt{code/verify\_spherical\_mean\_identity.py} checks
Lemma~\ref{lem:spherical} exactly in dimension three for
\(F(y)=|y|^4\), an arbitrary center, arbitrary scale, and a nontrivial radial
profile. It returns \texttt{symbolic\_difference=0} and \texttt{PASS}.
The script \texttt{code/verify\_second\_difference\_identity.py} performs the
corresponding exact one-dimensional check and also passes.
They were run with
\[
 \texttt{conda run --no-capture-output -n sandbox python code/<script>.py}.
\]

The two spectral scripts are exploratory only. They found no nonpositive
eigenvalue over the admissible amplitude range and thousands of irregular
horizons, but no computational observation is used in the proof.

\section*{Novelty check and limitations}

The run registry, solution, attempt, and proof-gap indexes were searched for
arXiv:2509.03468, the exact title, heterogeneous horizons, kernel triviality,
variable-bandwidth injectivity, Newtonian potentials, and spherical means.
The only indexed result was the dimension-one precursor now superseded by
this packet.

A bounded external search on July 23, 2026 used arXiv and exact combinations
of ``heterogeneous horizon'', ``nonlocal gradient kernel'', ``mildly
varying'', ``Newtonian potential'', and ``spherical mean maximum principle''.
It found the source preprint and unrelated nonlocal/variable-bandwidth work,
but no answer to Remark~4.18. The official arXiv record still lists only
version 1, submitted September 3, 2025. Novelty is plausible, not certified
by an exhaustive review.

The proof resolves precisely the source expectation under the source
hypotheses. It does not claim injectivity of \(\Qhet\) on every scalar
\(L^p\) function in dimension two: Lemma~\ref{lem:scalar} uses zero mass
there. The source auxiliary gradient has exactly that cancellation, which is
all Theorem~\ref{thm:main} requires.

\section*{Human review recommendation}

Review as a claimed full solution. The main checks are the singular endpoint
limits in Lemma~\ref{lem:spherical}, the strict quadratic perturbation in
Lemma~\ref{lem:max}, and the dimension-two boundary-component calculation in
Lemma~\ref{lem:moment}. A reviewer should also confirm the cited source proof
indeed supplies \(\widetilde w\in C_0^\infty(\overline\Omega)\) before mild
variation is invoked.

\begin{thebibliography}{9}

\bibitem{KS}
C. Kreisbeck and H. Schönberger,
\emph{Local boundary conditions in nonlocal hyperelasticity via heterogeneous
horizons},
arXiv:2509.03468v1, 2025.
\url{https://arxiv.org/abs/2509.03468}.

\bibitem{Evans}
L.C. Evans,
\emph{Partial Differential Equations},
2nd ed., Graduate Studies in Mathematics 19,
American Mathematical Society, 2010.

\end{thebibliography}

\end{document}
